\clearpage
\section{Prompts}
\label{app:prompts}

\subsection{MCP Server Annotation Prompt}
\label{app:server-category-annotation-prompt}

Below is the prompt for annotating MCP server categories.

\label{app:mcp-server-annotation-prompt}
\begin{minted}[fontsize=\footnotesize, linenos, breaklines, breakanywhere]{markdown}
## Task
Generate **Server Labels** to categorize the provided MCP Server based on its description and available tools.

## Objective
Analyze the provided MCP Server's description and available tools, then assign appropriate category labels that best describe its primary functionality and use cases.

## Guidelines

### Label Selection
- Analyze the MCP Server's core functionality and purpose
- Consider the types of tools it provides and the problems it solves
- Select labels that accurately represent the server's primary use cases
- Choose from predefined categories when applicable, but also consider custom labels for unique functionality

### Predefined Categories
Choose from these established categories when appropriate:
- **Web Search & Research**: Tools for searching the web, gathering information, academic research
- **Browser Automation**: Web scraping, automated browsing, page interaction
- **Memory Management**: Data storage, retrieval, knowledge bases, note-taking
- **Operating System**: File operations, system commands, process management
- **Data Analysis & Processing**: Analytics, data transformation, statistical analysis
- **Cryptocurrency & Blockchain**: Trading, wallet management, DeFi, blockchain interaction
- **Daily Productivity**: Task management, scheduling, personal organization
- **File Management**: File operations, document handling, storage management
- **Database Operations**: Data querying, database management, SQL operations
- **API Integration**: Third-party service integration, webhook handling
- **Communication Tools**: Messaging, email, notifications, social interaction
- **Development Tools**: Code analysis, debugging, version control, CI/CD
- **Security & Authentication**: Password management, encryption, access control
- **Cloud Services**: Cloud platform integration, serverless functions
- **AI/ML Tools**: Machine learning, model interaction, AI-powered features
- **Content Creation**: Writing, editing, media generation, publishing
- **Social Media**: Social platform integration, posting, analytics
- **Financial Services**: Banking, payments, financial data, accounting
- **E-commerce**: Shopping, product management, order processing
- **Gaming**: Game-related tools, entertainment, interactive features
- **Education**: Learning tools, course management, educational content
- **Health & Fitness**: Health monitoring, fitness tracking, medical tools
- **Travel & Maps**: Location services, travel planning, navigation
- **News & Media**: News aggregation, media consumption, journalism tools
- **Weather**: Weather data, forecasting, climate information
- **Time & Calendar**: Scheduling, time management, calendar integration

### Custom Labels
- If the server doesn't fit well into predefined categories, create a custom label
- Custom labels should be descriptive and specific to the server's unique functionality
- Use clear, concise terminology that would be useful for clustering and organization

### Output Requirements
- **Primary Label**: The main category that best describes the server (from predefined list or custom)
- **Secondary Labels**: Additional relevant categories (0-2 labels)
- **Custom Label**: A free-form descriptive label if the server has unique functionality not covered by predefined categories

## MCP Server Description
{MCP_SERVER_NAME}: {MCP_SERVER_DESCRIPTION}

Available Tools:
{TOOL_LIST}

## Output
Provide your response in the following XML format:

<response>
  <analysis>
    <!-- Briefly analyze the MCP Server's core functionality and the types of problems it solves based on its description and available tools. -->
  </analysis>
  <reasoning>
    <!-- Brief explanation of why these labels were chosen and how they represent the server's functionality -->
  </reasoning>
  <primary_label>
    <!-- The main category that best describes this server's primary functionality -->
  </primary_label>
  <secondary_labels>
    <!-- Additional relevant categories (0-2 labels), separated by commas if multiple -->
  </secondary_labels>
  <custom_label>
    <!-- A free-form descriptive label if the server has unique functionality not covered by predefined categories. Leave empty if not needed. -->
  </custom_label>
</response>
\end{minted}

\subsection{Task Generation Prompt}
\label{app:task-generation-prompt-for-single-server}

Below is an example of a task generation prompt for the single-server task synthesis. The prompt generates a question targeting \textbf{one tool}.

\begin{minted}[fontsize=\footnotesize, linenos, breaklines, breakanywhere]{markdown}
## Task
Generate a **Tool Use Question** based on the provided MCP Server and its tool descriptions.

## Objective
Analyze the provided MCP Server and its available tools, then create a realistic user question that would naturally require the use of one of these tools to solve.

## Guidelines

### Question Realism
- Create questions that represent real-world scenarios where users would need to interact with the MCP Server's tools
- The question should sound natural and authentic, as if asked by someone genuinely needing to accomplish a task
- Consider common use cases, problems, or workflows that would require the functionality provided by the MCP Server's tools

### Tool Selection
- Focus on **ONE specific tool** from the MCP Server that would be most appropriate to answer the question
- Choose tools based on the core functionality they provide and how they would solve real user problems
- Consider each tool's description and purpose when crafting the question

### Question Complexity
- Create questions that are clear and specific enough to warrant tool usage
- Avoid overly simple questions that could be answered without tools
- Include relevant context or constraints that make the tool usage necessary
- Do not contain the exact tool name in the question

### Output Format
Your response should include:
1. **Tool Analysis**: Briefly analyze the MCP Server's available tools and their main functionalities.
2. **Target Tool**: The specific tool name from the MCP Server that should be used to answer this question.
3. **Question**: A clear, realistic user question that requires tool usage.

## MCP Server Description
{MCP_SERVER_NAME}: {MCP_SERVER_DESCRIPTION}

Available Tools:
{TOOL_LIST}

## Output
Provide your response in the following XML format:

<response>
  <server_analysis>
    <!-- Briefly analyze the MCP Server's available tools and their main functionalities. -->
  </server_analysis>
  <target_tool>
    <!-- The specific tool name from the MCP Server that should be used to answer this question. -->
  </target_tool>
  <question>
    <!-- A clear, realistic user question that requires tool usage. -->
  </question>
</response>
\end{minted}

Below is an example of a task generation prompt for the single-server task synthesis. The prompt generates a question targeting \textbf{multiple tools}.

\begin{minted}[fontsize=\footnotesize, linenos, breaklines, breakanywhere]{markdown}
## Task
Generate a **Tool Use Question** based on the provided MCP Server and its tool descriptions.

## Objective
Analyze the provided MCP Server and its available tools, then create a realistic user question that would naturally require the use of **{NUM_TOOLS} tools** from this MCP Server to solve completely.

## Guidelines

### Question Realism
- Create questions that represent real-world scenarios where users would need to interact with the MCP Server's tools
- The question should sound natural and authentic, as if asked by someone genuinely needing to accomplish a task
- Consider common use cases, problems, or workflows that would require the functionality provided by the MCP Server's tools

### Tool Selection
- Focus on **{NUM_TOOLS} tools** from the MCP Server that would work together to answer the question
- The question should require a sequence or combination of tool calls to solve completely
- Choose tools based on how they complement each other and create a logical workflow
- Consider each tool's description and purpose when crafting the question that requires multiple steps

### Question Complexity
- Create questions that are complex enough to warrant using {NUM_TOOLS} tools
- The question should have multiple components or require several steps to solve
- Include relevant context or constraints that make the multi-tool usage necessary
- Do not contain the exact tool names in the question
- Ensure the question cannot be reasonably answered with just a single tool

### Output Format
Your response should include:
1. **Tool Analysis**: Briefly analyze the MCP Server's available tools and their main functionalities.
2. **Target Tools**: The specific tool names from the MCP Server that should be used together to answer this question, in the order they would likely be called.
3. **Question**: A clear, realistic user question that requires multiple tool usage.

## MCP Server Description
{MCP_SERVER_NAME}: {MCP_SERVER_DESCRIPTION}

Available Tools:
{TOOL_LIST}

## Output
Ensure your question requires exactly {NUM_TOOLS} tools to solve completely. Provide your response in the following XML format:

<response>
  <server_analysis>
    <!-- Briefly analyze the MCP Server's available tools and their main functionalities. -->
  </server_analysis>
  <target_tools>
    <!-- The specific tool names from the MCP Server that should be used together to answer this question, listed in order. e.g., <tool>create_twitter_post</tool> <tool>get_last_tweet</tool> -->
  </target_tools>
  <question>
    <!-- A clear, realistic user question that requires multiple tool usage. -->
  </question>
</response>
\end{minted}

Below is an example of a task generation prompt for the multi-server task synthesis. 

\begin{minted}[fontsize=\footnotesize, linenos, breaklines, breakanywhere]{markdown}
## Task
Generate a **Multi-Server Tool Use Question** based on the provided MCP Servers and their tool descriptions.

## Objective
Analyze the provided MCP Servers and their available tools, then create a realistic user question that would naturally require the use of **{NUM_TOOLS} tools from at least 2 different MCP servers** to solve completely.

## Guidelines

### Question Realism
- Create questions that represent real-world scenarios where users would need to interact with tools from multiple MCP Servers
- The question should sound natural and authentic, as if asked by someone genuinely needing to accomplish a complex task
- Consider workflows that span across different services/domains that would require multiple servers
- Think about how different MCP servers complement each other in real-world use cases

### Server and Tool Selection
- Use tools from **at least 2 different MCP servers** to answer the question
- Select **{NUM_TOOLS} tools total** that work together across multiple servers
- The question should require a sequence or combination of tool calls from different servers to solve completely
- Choose tools based on how they complement each other across different services/domains
- Consider each tool's description and purpose when crafting the cross-server workflow
- Ensure tools from different servers create a logical, interconnected workflow

### Question Complexity
- Create questions that are complex enough to warrant using {NUM_TOOLS} tools across multiple servers
- The question should have multiple components or require several steps that span different services
- Include relevant context or constraints that make the multi-server tool usage necessary
- Do not contain the exact tool names or server names in the question
- Ensure the question cannot be reasonably answered with tools from just a single server
- Create scenarios that naturally require different types of services working together

### Cross-Server Integration
- Think about how different servers' capabilities can be combined
- Consider data flow between different services (e.g., retrieving data from one service to use in another)
- Create realistic scenarios where multiple services need to work together
- Focus on complementary functionalities across different domains

### Output Format
Your response should include:
1. **Server Analysis**: Briefly analyze all MCP Servers and their available tools, focusing on how they can work together.
2. **Cross-Server Workflow**: Describe the workflow showing how tools from different servers will be used together.
3. **Target Tools**: The specific tool names from different MCP Servers that should be used together, in the order they would likely be called, with their server names.
4. **Question**: A clear, realistic user question that requires multi-server tool usage.

## Available MCP Servers

{SERVER_DESCRIPTIONS}

## Output
Ensure your question requires exactly {NUM_TOOLS} tools from at least 2 different servers to solve completely. Provide your response in the following XML format:

<response>
  <server_analysis>
    <!-- Briefly analyze all MCP Servers and their available tools, focusing on how they can work together across different domains/services. -->
  </server_analysis>
  <cross_server_workflow>
    <!-- Describe the workflow showing how tools from different servers will be used together to solve the question. -->
  </cross_server_workflow>
  <target_tools>
    <!-- The specific tool names from different MCP Servers that should be used together, listed in order with their server names. e.g., <tool server="Server1">search_posts</tool> <tool server="Server2">send_email</tool> -->
  </target_tools>
  <question>
    <!-- A clear, realistic user question that requires multi-server tool usage spanning different services/domains. -->
  </question>
</response> 
\end{minted}

Below is an example of a task generation prompt for the task synthesis for featured servers.

\begin{minted}[fontsize=\footnotesize, linenos, breaklines, breakanywhere]{markdown}
## Task
Generate a **Multi-Server Tool Use Question** based on featured MCP Servers and their tool descriptions.

## Objective
Brainstorm a compelling real-world scenario, then analyze the provided featured MCP Servers and their available tools to create a realistic user question that would naturally require the use of **{NUM_TOOLS} tools from at least 2 different MCP servers** to solve completely.

## Guidelines

### Scenario Brainstorming
- Think of realistic, specific scenarios where someone would need to use {NUM_TOOLS} different tools across multiple servers to accomplish a meaningful task
- Consider diverse real-world contexts such as:
  - Content creators managing their online presence across different platforms
  - Researchers gathering and analyzing information from multiple sources  
  - Developers building and deploying applications using different services
  - Business professionals managing projects and communications across platforms
  - Students working on complex assignments requiring multiple tools
  - Entrepreneurs launching new ventures using various services
- The scenario should be detailed and authentic, representing genuine use cases that span multiple services

### Question Realism
- Create questions that represent real-world scenarios where users would genuinely need tools from multiple MCP servers
- The question should sound natural and authentic, as if asked by someone with a specific goal
- Include relevant context, constraints, and details that make the question engaging
- Consider workflows that require multiple complementary tools working together across different services
- Think about how different servers support each other in real-world use cases

### Server and Tool Selection
- Use tools from **at least 2 different MCP servers** to answer the question
- Select **{NUM_TOOLS} tools total** that work together across multiple servers
- The question should require a sequence or combination of tool calls from different servers to solve completely
- Choose tools based on how they complement each other across different services/domains
- Consider each tool's description and purpose when crafting the cross-server workflow
- Ensure tools from different servers create a logical, interconnected workflow

### Question Complexity
- Create questions that are complex enough to warrant using {NUM_TOOLS} tools across multiple servers
- The question should have multiple components or require several steps that span different services
- Include relevant context or constraints that make the multi-server tool usage necessary
- Do not contain the exact tool names or server names in the question
- Ensure the question cannot be reasonably answered with tools from just a single server
- Create scenarios that naturally require different types of services working together

### Cross-Server Integration
- Think about how different servers' capabilities can be combined
- Consider data flow between different services (e.g., retrieving data from one service to use in another)
- Create realistic scenarios where multiple services need to work together
- Focus on complementary functionalities across different domains

### Output Format
Your response should include:
1. **Server Analysis**: Briefly analyze the featured MCP Servers and their available tools, focusing on how they can work together.
2. **Cross-Server Workflow**: Describe the workflow showing how tools from different servers will be used together.
3. **Target Tools**: The specific tool names from different MCP Servers that should be used together, in the order they would likely be called, with their server names.
4. **Question**: A clear, realistic user question that requires multi-server tool usage.

## Available Featured MCP Servers

{FEATURED_SERVER_DESCRIPTIONS}

## Output
Ensure your question requires exactly {NUM_TOOLS} tools from at least 2 different servers to solve completely. Provide your response in the following XML format:

<response>
  <server_analysis>
    <!-- Briefly analyze the featured MCP Servers and their available tools, focusing on how they can work together across different domains/services. -->
  </server_analysis>
  <cross_server_workflow>
    <!-- Describe the workflow showing how tools from different servers will be used together to solve the question. -->
  </cross_server_workflow>
  <target_tools>
    <!-- The specific tool names from different MCP Servers that should be used together, listed in order with their server names. e.g., <tool server="Server1">search_posts</tool> <tool server="Server2">send_email</tool> -->
  </target_tools>
  <question>
    <!-- A clear, realistic user question that requires multi-server tool usage spanning different services/domains. -->
  </question>
</response> 
\end{minted}

\subsection{Task Diversification Prompt}
\label{app:task-diversification-prompts}

The following prompt aims to add diversity to the given task by introducing new contexts and personas.

\begin{minted}[fontsize=\footnotesize, linenos, breaklines, breakanywhere]{markdown}
## Task
Generate **augmented variations** of a given question that maintain the same target tool(s) usage and complexity level but apply them across different contexts and scenarios.

## Objective
Take an existing question and its associated target tool(s), then create multiple variations that:
- Use the same target tool(s) to achieve the core goal
- Maintain the exact same tool usage order and final outcome
- Apply the question to completely different contexts, scenarios, or domains
- Keep the same level of complexity and constraints as the original
- Demonstrate how the same tool usage pattern applies across diverse real-world scenarios

## Guidelines
- Translate the question to distinctly different domains, user personas, or situational contexts while preserving its original complexity level.
- Keep the tool usage sequence and final outcome identical across all variations.
- Ensure each variation feels like a realistic scenario in its new context and remains solvable with the same tool operations.
- Ensure the question does not contain any tool names or explicit references to the target tools.

## Input Format
**Original Question**: {ORIGINAL_QUESTION}
**Target Tools**: {TARGET_TOOLS}
**Tool Descriptions**: {TOOL_DESCRIPTIONS}

## Output Requirements
Generate **{VARIATIONS_COUNT} augmented variations** of the original question. Each variation should:
1. Maintain the same core goal that requires the target tool(s)
2. Use the exact same tool(s) in the same order with the same final outcome
3. Apply to a completely different context, scenario, or domain
4. Keep the same complexity level and constraints as the original
5. Feel like a natural, real-world scenario from a different setting
6. Be meaningfully different from the original and other variations in terms of context only
7. Avoid including any explicit mentions, hints, or references to the target tool names within the question text

## Output
Provide your response in the following XML format:

<response>
  <analysis>
    <!-- Briefly analyze the original question and target tool(s) to understand the core goal, tool usage pattern, complexity level, and expected outcome, then identify how this can be applied across different domains while maintaining operational consistency -->
  </analysis>
  <variations>
    <!-- Generate {VARIATIONS_COUNT} variations, each with <variation_X>, <context>, and <question> tags -->
    <variation_1>
      <context>
        <!-- Brief description of the new domain/scenario introduced -->
      </context>
      <question>
        <!-- The augmented question that maintains the same target tool(s) usage order, complexity, and outcome but in a different context -->
      </question>
    </variation_1>
    <!-- Continue with variation_2, variation_3, etc. as needed based on number of variations -->
  </variations>
</response>
\end{minted}

The prompt below is designed to enhance task complexity through the introduction of additional constraints.

\begin{minted}[fontsize=\footnotesize, linenos, breaklines, breakanywhere]{markdown}
## Task
Generate **augmented variations** of a given question that maintain the same target tool(s) usage and context but significantly increase the complexity and constraints required to solve the problem.

## Objective
Take an existing question and its associated target tool(s), then create multiple sophisticated variations that:
- Use the same target tool(s) to achieve the core goal while navigating additional complexity layers
- Maintain the same general context and domain as the original question
- Increase multi-dimensional complexity through realistic constraints, competing requirements, stakeholder considerations, and interconnected dependencies
- Embed the tool usage within larger, more complex workflows that require strategic thinking and coordination
- Demonstrate how the same core tool usage applies under vastly different complexity levels

## Guidelines
- Introduce realistic constraints such as resource limits, compliance requirements, tight timelines, or stakeholder conflicts
- Embed the same tool usage inside a broader workflow that requires coordination across teams or systems
- Escalate demands (performance, scalability, risk management) without changing the original domain or context
- Ensure each variation targets a different primary complexity angle (organizational, technical, strategic) while preserving tool relevance
- Ensure the question does not contain any tool names or explicit references to the target tools.

## Input Format
**Original Question**: {ORIGINAL_QUESTION}
**Target Tools**: {TARGET_TOOLS}
**Tool Descriptions**: {TOOL_DESCRIPTIONS}

## Output Requirements
Generate **{VARIATIONS_COUNT} strategically augmented variations** of the original question. Each variation should:
1. Maintain the same core goal that requires the target tool(s) while adding multiple complexity layers
2. Keep the same general context and domain as the original question
3. Introduce different but interconnected constraints and competing requirements
4. Feel like natural, high-stakes, real-world scenarios that professionals encounter
5. Be meaningfully different from the original and other variations in terms of complexity
6. Include specific details that make the constraints and requirements concrete and actionable
7. **Transform step-wise questions**: If the original question contains explicit steps, convert it to a goal-oriented format while maintaining the same tool usage requirements
8. Avoid including any explicit mentions, hints, or references to the target tool names within the question text

## Output
Provide your response in the following XML format:

<response>
  <analysis>
    <!-- Analyze the original question and target tool(s) to understand the core goal, current complexity level, and identify multiple complexity dimensions that can be naturally introduced while maintaining tool relevance and solution feasibility -->
  </analysis>
  <variations>
    <!-- Generate {VARIATIONS_COUNT} variations, each with <variation_X>, <constraints>, and <question> tags -->
    <variation_1>
      <constraints>
        <!-- Specific organizational, stakeholder, or coordination constraints that add realistic complexity -->
      </constraints>
      <question>
        <!-- The complex, organizationally-focused question that maintains the same target tool(s) usage within a more sophisticated workflow -->
      </question>
    </variation_1>
    <!-- Continue with variation_2, variation_3, etc. as needed based on number of variations -->
  </variations>
</response>

\end{minted}

\subsection{Task Quality Annotation Prompt}
\label{app:task-annotation-prompt-for-single-server}
\begin{minted}[fontsize=\footnotesize, linenos, breaklines, breakanywhere]{markdown}
## Task
Conduct a **Question Quality Assessment** of a tool use question across six key dimensions to ensure it meets high standards for realistic tool usage scenarios.

## Objective
Analyze the provided tool use question and assess its quality across six primary dimensions:
1. **Tool Selection Difficulty** - How challenging it is to determine which tools to use giving all available tools
2. **Tool Selection Uniqueness** - How unique and necessary the selected tools are for this specific task giving all available tools
3. **Question Quality** - Overall clarity, specificity, and effectiveness 
4. **Scenario Realism** - How authentic and believable the scenario is
5. **Verifiable** - How easy it is to verify the correctness of the final model answer
6. **Stability** - How stable the answer will be when requested under different time and geolocation

## Assessment Criteria

### 1. Tool Selection Difficulty
**What to Evaluate**: How difficult it would be for a user to determine which specific tools are needed to solve this question.

**Rating Guidelines**:
- **very easy**: Question explicitly mentions tool names or makes tool selection obvious
- **easy**: Tool selection is straightforward with clear indicators
- **medium**: Requires some reasoning but tool needs are fairly apparent  
- **hard**: Requires careful analysis to determine appropriate tools
- **very hard**: Requires extensive expertise and deep reasoning to identify the correct tools

### 2. Tool Selection Uniqueness
**What to Evaluate**: How unique and necessary the selected tools are for accomplishing this specific task, and whether the task can only be completed with these tools in the specified sequence.

**Rating Guidelines**:
- **not unique**: Many alternative tool combinations could accomplish the same task equally well
- **somewhat unique**: Some alternative approaches exist, but selected tools offer advantages
- **moderately unique**: Selected tools are well-suited, with limited alternative approaches
- **quite unique**: Selected tools are particularly well-matched to the task requirements
- **highly unique**: Task can only be accomplished effectively with these specific tools in this sequence

### 3. Question Quality
**What to Evaluate**: Overall quality, clarity, and effectiveness of the question as a realistic user query.

**Rating Guidelines**:
- **very poor**: Unclear, ambiguous, or poorly constructed question
- **poor**: Some clarity issues, missing important context
- **average**: Clear and understandable, but could be more specific or engaging
- **good**: Well-constructed, clear, specific, and realistic
- **excellent**: Exceptionally clear, detailed, engaging, and professionally written

### 4. Scenario Realism
**What to Evaluate**: How authentic, believable, and true-to-life the described scenario is.

**Rating Guidelines**:
- **unrealistic**: Artificial, contrived, or implausible scenario
- **somewhat unrealistic**: Some realistic elements but feels forced or unlikely
- **moderately realistic**: Believable scenario with minor authenticity issues
- **realistic**: Authentic scenario that represents genuine use cases
- **highly realistic**: Completely natural, authentic scenario indistinguishable from real user needs

### 5. Verifiable
**What to Evaluate**: How easy it is to verify the correctness of the final model answer.

**Rating Guidelines**:
- **hard to verify**: Fully free-form answer that requires extensive human judgment
- **somewhat hard**: Mostly subjective answer with some verifiable elements
- **moderately verifiable**: Short sentence that can be verified by LLM comparison
- **mostly verifiable**: Answer with clear, objective components and some subjective elements
- **easy to verify**: Answer can be verified by simple rules, exact matches, or clear success criteria

### 6. Stability (1-5 Scale)
**What to Evaluate**: How stable and consistent the answer will be when the question is asked under different environmental conditions and system contexts. Consider factors like temporal dependency, geographical variations, operating system differences, network environments, and software version variations.

**Rating Guidelines**:
- **highly unstable**: Answer changes significantly across different conditions (real-time data, location-specific, system-dependent)
- **somewhat unstable**: Answer may vary moderately based on environmental or system factors
- **moderately stable**: Answer mostly consistent with minor variations due to context
- **mostly stable**: Answer remains largely consistent across different conditions
- **highly stable**: Answer is completely independent of environmental and system factors

## Question Analysis

### All Available Tools```
{ALL_SERVER_AND_TOOL_INFORMATION}
```

### Question Content
```
{QUESTION_CONTENT}
```

### Intended Tool for This Question
```
{INTENDED_TOOL}
```

## Output Requirements

Provide analysis with detailed reasoning BEFORE scores for each of the six metrics.

## Output
Provide your response in the following XML format:

<response>
  <tool_selection_difficulty>
    <reasoning>
      <!-- Detailed explanation including ambiguity level, domain knowledge required, and alternative solutions giving all available tools -->
    </reasoning>
    <rating><!-- Rating: very easy, easy, medium, hard, very hard --></rating>
  </tool_selection_difficulty>
  
  <tool_selection_uniqueness>
    <reasoning>
      <!-- Detailed explanation of tool necessity, sequential dependencies, and alternative tool viability giving all available tools -->
    </reasoning>
    <rating><!-- Rating: not unique, somewhat unique, moderately unique, quite unique, highly unique --></rating>
  </tool_selection_uniqueness>
  
  <question_quality>
    <reasoning>
      <!-- Detailed explanation covering linguistic quality, information architecture, and actionability -->
    </reasoning>
    <rating><!-- Rating: very poor, poor, average, good, excellent --></rating>
  </question_quality>
  
  <scenario_realism>
    <reasoning>
      <!-- Detailed explanation of industry authenticity, workflow accuracy, and stakeholder behavior -->
    </reasoning>
    <rating><!-- Rating: unrealistic, somewhat unrealistic, moderately realistic, realistic, highly realistic --></rating>
  </scenario_realism>
  
  <verifiable>
    <reasoning>
      <!-- Detailed explanation of answer format, objective criteria, and ground truth availability -->
    </reasoning>
    <rating><!-- Rating: hard to verify, somewhat hard, moderately verifiable, mostly verifiable, easy to verify --></rating>
  </verifiable>
  
  <stability>
    <reasoning>
      <!-- Detailed explanation of temporal/geographical/system dependencies and environmental factors -->
    </reasoning>
    <rating><!-- Rating: highly unstable, somewhat unstable, moderately stable, mostly stable, highly stable --></rating>
  </stability>
</response> 

\end{minted}

\subsection{Trajectory Annotation Prompt}
\label{app:trajectory-annotation-prompt-for-single-server}
\begin{minted}[fontsize=\footnotesize, linenos, breaklines, breakanywhere]{markdown}
## Task
Conduct a **Response Quality Assessment** of a tool-use conversation across two LLM-scored dimensions, with a third dimension computed automatically outside the LLM.

## Objective
Analyze the provided conversation and assess its response quality across two primary dimensions scored by the LLM, while reserving an additional tool-call accuracy dimension for automated scoring:
1. Completeness - Whether the assistant fully accomplished the user's request end-to-end
2. Conciseness - Whether the assistant solved the task using the minimum necessary steps and verbosity

## Assessment Criteria

### 1. Completeness
**What to Evaluate**: Did the assistant fully satisfy the user's goal given the conversation context? Consider whether the assistant:
- Executed all required steps end-to-end (including saving/exporting/downloading where applicable)
- Provided the final deliverable or a working alternative when blocked (e.g., tool failure with a usable fallback)
- Included essential confirmations, paths, or instructions to achieve the outcome
- Avoided missing key requirements or leaving the user with unresolved gaps

**Rating Guidelines**:
- very incomplete: Major requirements missing; no usable outcome
- incomplete: Some key requirements missing; outcome is not directly usable
- partially complete: Core steps attempted; outcome usable only with user effort or missing minor requirements
- mostly complete: Meets most requirements; small omissions or minor issues remain
- fully complete: All requirements met with a usable outcome delivered

### 2. Conciseness
**What to Evaluate**: Did the assistant achieve the goal with minimal redundancy and steps? Consider whether the assistant:
- Avoided repetitive or unnecessary explanations/tool calls
- Used the minimal set of steps/tools to complete the task
- Kept language concise while preserving clarity

**Rating Guidelines**:
- very redundant: Excessive repetition or unnecessary steps/tool calls
- redundant: Noticeable verbosity or extra steps beyond what's needed
- average: Reasonably concise with minor extraneous content
- concise: Efficient and to the point with minimal overhead
- very concise: Maximally efficient while clear and complete

## Response Analysis

### Question Content
```
{QUESTION_CONTENT}
```

### Intended Tool for This Question
```
{INTENDED_TOOL}
```

### Conversation History
```
{CONVERSATION_HISTORY}
```

## Output Requirements
- Provide detailed reasoning BEFORE ratings for Completeness and Conciseness
- Do NOT score Tool Call Accuracy; include placeholders only

## Output
Provide your response in the following XML format:

<response>
  <completeness>
    <reasoning>
      <!-- Evaluate if the assistant delivered an end-to-end usable outcome, addressed all requirements, handled tool failures with alternatives, and provided necessary confirmations/paths. -->
    </reasoning>
    <rating><!-- Rating: very incomplete, incomplete, partially complete, mostly complete, fully complete --></rating>
  </completeness>

  <conciseness>
    <reasoning>
      <!-- Evaluate if the assistant minimized redundant steps/explanations, avoided unnecessary tool calls, and kept messaging efficient while clear. -->
    </reasoning>
    <rating><!-- Rating: very redundant, redundant, average, concise, very concise --></rating>
  </conciseness>
</response>
\end{minted}
